\section{A Protocol That Could Settle It: the Canine Ladder Benchmark}
\label{sec:protocol}

The complaint running through Sections~\ref{sec:products}--\ref{sec:eval} is
methodological, not skeptical: this field reports accuracies against labels
that no independent process can check. A survey that diagnoses a measurement
problem owes the field a measurement, so this section specifies one. The
design rule is single and severe. \textbf{Each rung may be scored only against
an error signal the model cannot see and the experimenter cannot leak.}
Everything below follows from that rule, and each rung gets a task, the ground
truth that makes the task checkable, a primary metric, a baseline the system
must beat, and the observation that would show the rung was not actually
reached.

\subsection{The five tasks}

\textbf{Rung 1: signal and sender.} Separate vocalisation from non-vocal sound,
then identify the individual. The error signal is objective: session metadata
and a video-verified sender, the trick that made the individual-identification
studies convincing~\cite{molnar2006discriminate}. Score equal-error rate on
splits taken per dog and per recording session, so that microphone and room
cannot carry identity. The baseline that must be beaten is a human
speaker-verification embedding transferred without fine-tuning, because that
tests whether canine vocal anatomy contributes anything a human-speaker model
does not already supply~\cite{snyder2018xvector}. A system has not reached
rung 1 if its accuracy survives a session shuffle: that means it found the
recording equipment, not the dog.

\textbf{Rung 2: affect.} Classify the animal's affective state over a
pre-registered taxonomy. Human consensus alone is not ground truth here --- it
is a proxy, and treating it as a readout is the field's most common
substitution. A scored submission needs at least two anchors fixed before any
model sees the data: a manipulation the experimenter controls and records in
advance (isolation, food barrier, territorial stranger, play; the design
Pongr\'acz and Farag\'o used~\cite{pongracz2005barks,farago2017growls}), and an
animal-side readout that does not depend on anyone's judgement --- heart-rate
variability, latency to approach or withdraw, latency to settle. Human ratings
are admitted as a third anchor only if collected blind to the manipulation and
never used for training. Report balanced accuracy together with calibration,
because an over-confident affect classifier in a consumer app is worse than a
honest one that abstains~\cite{guo2017calibration}, and report an abstention
curve. The baseline that must be beaten is an arousal-only regressor built from
energy and fundamental-frequency statistics. If it ties the model, the model
measured arousal, not affect, and the paper must say so in those words.

\textbf{Rung 3: intent and reference.} Predict what the animal is acting on.
The only clean error signal is the animal's own behaviour inside a forced-choice
design whose alternatives are physically constrained: an object-choice task with
the olfactory route eliminated, which is exactly the control that made the word
- comprehension results defensible~\cite{kaminski2004rico,pilley2011chaser}, or a
button array recorded while the handler is blind to the target. Score accuracy
against a permutation baseline over the observed sequences, as the
combination-analysis study did~\cite{bastos2024scirep}, and require
generalisation across dogs. The rung is not reached if performance collapses
when the human leaves the room; that collapse is the Clever Hans signature, and
a benchmark that never tests it will keep mistaking cueing for competence.
The alternative is already standard practice in animal cognition: cross-modal
matching experiments score what the animal does --- which of two people a
horse looks at after hearing one of them speak --- rather than what a human
judge says the recording
means~\cite{proops2009horses,proops2012familiar}. That swap is the difference
between a measurement and a performance, and it is the single most important
change this benchmark asks of canine work.

\textbf{Rung 4: referential semantics, tested as composition.} Map multi-element signals or sequences onto
multi-element states. The error signal is a closed inventory of objects and
actions with combinations held out of training, and the metric is accuracy on
unseen combinations against a co-occurrence factorisation that has seen the same
elements but not the same pairings. A model that matches the factorisation has
learned frequencies, not structure. This is the rung where
Figure~\ref{fig:evidence} shows the validated column empty, so a single
surviving entry here would be the most informative result the benchmark could
produce in its first year.

\textbf{Rung 5: open-ended dialogue.} There is no protocol that can certify this
rung, and the survey will not pretend otherwise: certifying it would require an
independent record of an intended sentence, and no such record exists for a
species without observable language. What can be built is the other half of
science --- a test that can reject. Run a behavioural-equivalence trial: give
naive human participants a rendered sentence and a control sentence matched for
affect, polarity and length, then measure what the participants subsequently do
for the dog. If the two conditions are indistinguishable, the rendering carries
no information beyond rung 2 and the word ``translation'' is refuted for that
system. The test is one-sided by construction. It can close claims down; it
cannot open a rung up, and a benchmark report should state which of the two it
is doing.

\begin{table}[htbp]
\centering
\caption{The Canine Ladder Benchmark: one task per rung, with the error signal
that makes it checkable and the baseline a submission must beat. The final
column is the observation that shows the rung was not actually reached.}
\label{tab:protocol}
\footnotesize
\setlength{\tabcolsep}{4pt}
\begin{tabular}{@{}p{1.5cm}p{3.1cm}p{3.5cm}p{3.4cm}p{3.4cm}@{}}
\toprule
\textbf{Rung} & \textbf{Task} & \textbf{Error signal} & \textbf{Must beat} & \textbf{Not reached if} \\
\midrule
1 Signal & vocal vs.\ non-vocal; identify the dog & session metadata, video-verified sender & transferred human speaker-verification embedding & accuracy survives a session shuffle \\
2 Affect & $n$-class state over a pre-registered taxonomy & controlled pre-recorded manipulation $+$ animal-side physiological or latency readout & arousal-only energy/$f_0$ regressor & the arousal baseline ties the model \\
3 Intent & what the animal is acting on & forced-choice behaviour with alternatives physically constrained; handler blind & permutation baseline over sequences & performance collapses when the human leaves \\
4 Reference & multi-element signal $\rightarrow$ multi-element state & held-out combinations from a closed inventory & co-occurrence factorisation & factorisation matches the model \\
5 Dialogue & render an intended sentence & none available; use a rejection test & --- & see the text: this rung is scored by falsification only \\
\bottomrule
\end{tabular}
\end{table}

\subsection{Rules that apply to every rung}

Four cross-cutting requirements do most of the work, and none of them is
novel --- the point is that this field has not adopted them. Splits are per dog
and per session, never per clip; a benchmark whose test set contains another
recording of the same animal in the same room will be solved by the room.
Taxonomies are pre-registered, so that a class set cannot be edited until the
confusion matrix looks good. Reporting follows the artifact discipline the
ML side of this survey already has vocabulary for: datasets ship with
datasheets~\cite{gebru2021datasheets}, models with model
cards~\cite{mitchell2019modelcards}, and a claim about capability is written at
the rung the evidence reaches, which is the discipline whose absence produced
the sentences in Table~\ref{tab:products}. And abstention is legal and scored:
a system that says ``insufficient signal'' on half of a home recording and is
right about the other half is more useful than one that always guesses.

A submission is then judged on a short checklist, and the checklist is the
contribution as much as the tasks are: per-dog split declared; error signal
named and independent of the annotator; arousal baseline reported; human
raters blind; handler absent or blinded; calibration reported; abstention
reported; negative controls included; taxonomy declared before training;
rung stated in the abstract. Ten items, answerable in a page, and every one of
them corresponds to a way a current published or marketed claim has already
failed.

\subsection{What adoption would buy}

A shared benchmark changes the incentive structure, which is what the
evaluation crisis actually is. Today the cheapest way to look good is to move
the label, since labels are cheap and error signals are not; under this
protocol the only way up the ladder is to acquire a better error signal, which
means better experiments with dogs. The first version does not require new
animal work: the existing falsifiable datasets --- bark-context recordings,
growl corpora, the button-array logs, the object-choice archives --- already
cover rungs 1--3, and a community release could stand them up in months. The
hard part is institutional rather than technical, and it is the recommendation
of this survey: a small consortium of the psychophysics groups
that produced those corpora, a cross-species decoding lab, and one independent
auditor should own the split, the taxonomies and the leaderboard, and vendor
systems should be invited to submit to it rather than report numbers about
themselves. Whether any of that happens is the real test of whether the field
wants to know what it has built.
